% TOP_PROBLEM: {"id": 2302039, "title": "Research Problems in Function Theory — Problem 2.39", "category_id": 8, "category": {"id": 8, "name": "analysis", "display_name": "Analysis", "description": "Limits, continuity, calculus, and function theory.", "slug": "analysis", "order_index": 8, "created_at": "2026-07-31T15:26:25.671Z"}, "status": "partially_solved", "problem_number": "AMR-022-2039", "set_id": 13}
% FIRST_POSTED: 2026-10-02
% ENTRY_KIND: statement_only
% UnsolvedMath ID 2302039; source catalogue code AMR-022-2039
% !TeX program = lualatex
% Definition and sources only; this file is not an LLM solution attempt.
\documentclass[11pt,a4paper]{article}
\usepackage[margin=25mm]{geometry}
\usepackage{fontspec}
\setmainfont{lmroman10-regular.otf}[BoldFont=lmroman10-bold.otf,
 ItalicFont=lmroman10-italic.otf,BoldItalicFont=lmroman10-bolditalic.otf]
\usepackage{amsmath,amssymb,mathtools}
\usepackage{unicode-math}
\setmathfont{latinmodern-math.otf}
\AtBeginDocument{\let\setminus\smallsetminus}
\usepackage{xurl,hyperref}
\hypersetup{colorlinks=true,urlcolor=blue}
\setcounter{secnumdepth}{0}
\setlength{\parindent}{0pt}
\setlength{\parskip}{5pt}
\setlength{\emergencystretch}{3em}
\begin{document}
\section*{Research Problems in Function Theory — Problem 2.39}\label{2302039}
{\small UnsolvedMath ID 2302039 · AMR-022-2039 · Analysis · AMR Open Problem Lists}
\subsection{Definitions and mathematical statement}
For a transcendental entire function $f$, let
\[
T(r,f)=\frac1{2\pi}\int_0^{2\pi}\log^+|f(re^{i\theta})|\,d\theta,
\qquad m_0(r,f)=\min_{|z|=r}|f(z)|,
\]
where $\log^+x=\max(0,\log x)$. Its lower order is
$\lambda(f)=\liminf_{r\to\infty}\log T(r,f)/\log r$.
For each prescribed finite $\lambda\ge1$, determine the sharp constant
\[
D(\lambda)=\inf_{\lambda(f)=\lambda}
\limsup_{r\to\infty}\frac{\log m_0(r,f)}{T(r,f)}.
\]
This is the largest number for which every entire function of that lower order satisfies the corresponding lower bound for the upper limit. Zeros on a circle contribute $\log m_0=-\infty$, and do not remove that circle from the definition.
The smaller-order range is supplied as context by the source:
$D(\lambda)=\pi\lambda\cot(\pi\lambda)$ for $0\le\lambda<1/2$, with the limiting value at zero, and $D(\lambda)=\pi\lambda\cos(\pi\lambda)$ for $1/2\le\lambda<1$.
The task is the unresolved sharp lower-order range at and above one; a lower bound without matching extremal examples does not determine $D(\lambda)$.
\subsection{Short English statement}
Find the sharp lower bound for the normalized logarithmic minimum modulus of an entire function with prescribed lower order at least one.
\subsection{Sources}
\begin{itemize}
\item[S1] ULAM.AI and the UnsolvedMath Contributors, supplied problems.json, record 2302039 (AMR-022-2039), AMR Open Problem Lists. Catalogue identity and source transcription; CC BY 4.0.
\url{https://huggingface.co/datasets/ulamai/UnsolvedMath}.
\item[S2] Hayman - Research Problems in Function Theory (2018), Problem 2.39. Original source indexed by AMR-022-2039.
\url{https://arxiv.org/abs/1809.07200}.
\end{itemize}
\end{document}
